\chapter{Usage instructions}

\section{Test build}
To run unit tests:
\begin{lstlisting}[style=cppstyle]
make test
\end{lstlisting}
This command invokes the script and produces the test executable in the \texttt{test/build} folder.

To clean:
\begin{lstlisting}[style=cppstyle]
make test-clean
\end{lstlisting}
Cleaning is recommended when switching toolchains or updating dependencies.

\section{Arduino compilation}
To compile the sketch via the script:
\begin{lstlisting}[style=cppstyle]
./test/run_tests.sh --with-arduino --board arduino:avr:uno
\end{lstlisting}
This mode runs unit tests first and then compiles the sketch, reducing the risk of uploading unverified firmware.

To upload to a real board:
\begin{lstlisting}[style=cppstyle]
./test/run_tests.sh --upload --board arduino:avr:uno --port /dev/ttyACM0
\end{lstlisting}
Before uploading, verify that the serial port is correct.

\section{Wokwi simulation}
The simulation wiring is defined in \texttt{diagram.json}. Open the project in Wokwi and start the simulation to observe the overall behavior. The simulation is useful for validating the flow without physical hardware.
